% Part: counterfactuals
% Chapter: minimal-change-semantics
% Section: introduction

\documentclass[../../../include/open-logic-section]{subfiles}

\begin{document}

\olfileid{cnt}{min}{int}

\olsection{مقدمه}

استالنیکر و لوئیس گزارش‌هایی از شرطی‌های خلاف واقع، مانند «اگر کبریت
کشیده می‌شد، روشن می‌شد»، پیشنهاد کردند. گزارش‌های آنان پیشنهادهایی
دربارهٔ چگونگی فهم درست شرایط صدق چنین جمله‌هایی بود. اندیشهٔ مشترک هر
دو پیشنهاد این است: برای ارزیابی صدق یک شرطی خلاف واقع، باید جهان‌های
ممکنی را در نظر بگیریم که با کمینهٔ تفاوت از جهان بالفعل، مقدم را صادق
می‌سازند. اگر تالی در این جهان‌های ممکن صادق باشد، شرطی خلاف واقع صادق
است. برای نمونه، فرض کنید کبریت و کبریت‌دانی در دست دارم. در جهان بالفعل
تنها به آن‌ها نگاه می‌کنم و می‌اندیشم اگر کبریت را می‌کشیدم چه رخ
می‌داد. کمینهٔ تغییر نسبت به جهان بالفعل که در آن کبریت را می‌کشم این
است که تصمیم به عمل می‌گیرم و کبریت را می‌کشم. تغییر کمینه است، زیرا
هیچ چیز دیگری عوض نمی‌شود: هم‌زمان به هوا نمی‌پرم، کشیدن کبریت موهایم
را نیز آتش نمی‌زند، ناگهان همهٔ نیروی انگشتانم را از دست نمی‌دهم، در
کمین با تفنگ آب‌پاش هم‌زمان خیس نمی‌شوم، و مانند آن. در آن امکان بدیل،
کبریت روشن می‌شود. پس صادق است که اگر کبریت را می‌کشیدم، روشن می‌شد.

می‌توان این گزارش شهودی را با معناشناسی صوری برای منطق‌های شرطی خلاف
واقع همراه کرد. لوئیس نماد «$\boxright$» و استالنیکر نماد~«$>$» را برای
شرطی خلاف واقع معرفی کردند. ما از~$\cif$ استفاده می‌کنیم و آن را به‌صورت
رابطی دوتایی به منطق گزاره‌ای می‌افزاییم. بنابراین افزون بر
!!{formula}های صورت~$!A \lif !B$، !!{formula}های صورت~$!A \cif !B$ را
نیز داریم. معناشناسی صوری، مانند معناشناسی رابطه‌ای منطق موجهات، بر
مدل‌هایی مبتنی است که در آن‌ها !!{formula}ها در جهان‌ها ارزیابی
می‌شوند؛ و شرط ارضای تعریف‌کنندهٔ~$\mSat{M}{!A \cif !B}[w]$ برحسب
$\mSat{M}{!A}[w']$ و~$\mSat{M}{!B}[w']$ برای برخی جهان‌های (دیگر)
$w'$ داده می‌شود. کدام~$w'$؟ به‌طور شهودی، نزدیک‌ترین جهان (یا جهان‌ها)
به~$w$ که در آن~$\mSat{M}{!A}[w']$. این امر مستلزم آن است که رابطه‌ای
از «نزدیکی» نیز در مدل گنجانده شود.

لوئیس شیوه‌ای روشنگر برای بازنمایی گرافیکی وضعیت‌های خلاف واقع معرفی
کرد. هر جهان ممکن در مرکز مجموعه‌ای از کره‌های تودرتو قرار دارد که
جهان‌های دیگر را دربرمی‌گیرند---این کره‌ها را به‌صورت دایره‌های هم‌مرکز
می‌کشیم. جهان‌های میان دو کره به یک اندازه به جهان مرکزی نزدیک‌اند؛
جهان‌های درون کره‌ای داخلی‌تر نزدیک‌تر و جهان‌های درون کره‌ای بیرونی‌تر
دورترند.
\begin{center}
\begin{tikzpicture}[scale=.7]\small
  \spheresystem{5}
  \draw (0,0) node {$w$};
  \propositionintersect{0}{4}{40}{3.5}
  {\tikzset{proposition/.append={smooth,tension=1.4}}
  \proposition[shift={(1.6,-1)}]{90}{1}{30}{4}}
  \path (1.5,3) node[below] {$\formula{B}$};
  \path (3,0) node {$\formula{A}$};
\end{tikzpicture}
\end{center}
نزدیک‌ترین جهان‌های~$!A$ همان جهان‌های~$w'$ای هستند که~$!A$ در آن‌ها
ارضا می‌شود و در کوچک‌ترین کره پیرامون جهان مرکزی~$w$ قرار دارند (ناحیهٔ
خاکستری). به‌طور شهودی، $!A \cif !B$ در~$w$ ارضا می‌شود اگر~$!B$ در
همهٔ نزدیک‌ترین جهان‌های~$!A$ صادق باشد.

\end{document}
